% Secciones 4--7 del delta gravitatorio íntegro.
\section{Hexad--octad incidence, Barbero--Immirzi parameter, and area spectrum}
\label{sec:l9s102:barbero-area}

The areal derivation receives as antecedents the incidence inclusion
\(6\to8\), the TPK cardinals \(90/120\), the normalized defect
\(-1/12\), and the spectral channels \(q_\pm\). The
hexad--octad functional determines the Barbero--Immirzi parameter and the
normalization of the area operator. This residence contains the HMT derivation
and the spectrum; the Palatini--Cartan--Holst realization occupies a subsequent
residence in Dimensional Mechanics.

Let \(H\subset O\), with \(|H|=6\) and \(|O|=8\), and preserve the marked
pair of the hexad in the inclusion
\begin{equation}
 \iota_{6,8}:\mathcal F_6=H\times\binom H2
 \lhook\joinrel\longrightarrow
 \mathcal F_8=O\times\binom H2.
 \label{eq:l9s102:barbero-inclusion-incidencia}
\end{equation}
The internal count gives
\begin{equation}
 |\mathcal F_6|=6\binom62=90,
 \qquad |\mathcal F_8|=8\binom62=120,
 \qquad
 \frac{90-120}{24\binom62}=-\frac1{12}.
 \label{eq:l9s102:barbero-conteos}
\end{equation}
Once \(\pi_{\rm HMT}\), the angle \(A\) and the
unique counterangle \(C^*\) have been published, the two channels are
\begin{equation}
 q_\pm=\exp\!\left[-\frac{\pi_{\rm HMT}}{180}(A\pm C^*)\right],
 \qquad 0<q_+<q_-<1,
 \label{eq:l9s102:barbero-qpm}
\end{equation}
and their oriented-return series are
\begin{equation}
 S_{90}(q)=\frac{q^{90}}{1-q^{270}},\qquad
 S_{120}(q)=\frac{q^{120}}{1-q^{360}},\qquad
 \Delta S_m=S_m(q_-)-S_m(q_+).
 \label{eq:l9s102:barbero-series}
\end{equation}

The causal order of the coefficients is fixed in the temporal register of the TPK,
before evaluating the pole coefficient. Let \(C_{12}=\mathbb Z/12\mathbb Z\) be the dodecaphase
wheel and let \(\partial_{\rm ret}\) be its unique oriented return seam.
In the free module of incidence events, the fundamental turn is
\begin{equation}
 c_{12}^{+}
 =\sum_{j\in C_{12}}[j,\mathcal F_6]
  -[\partial_{\rm ret},\mathcal F_8].
 \label{eq:l9s102:barbero-cadena-dodecafasica}
\end{equation}
The sign of the second term is the boundary sign. The bidirectional
involution reverses the complete orientation,
\(c_{12}^{-}=-c_{12}^{+}\), without modifying the absolute multiplicities of
the two incidence types. Define the return reader on the
generators by
\begin{equation}
 \mathscr R([j,\mathcal F_6])=S_{90},
 \qquad
 \mathscr R([\partial_{\rm ret},\mathcal F_8])=S_{120}.
 \label{eq:l9s102:barbero-lector-retorno}
\end{equation}
Phase covariance makes the first image constant over the twelve
sectors; the second acts only at the enlarged seam. By linearity,
\begin{equation}
 \boxed{\mathscr R(c_{12}^{+})=12S_{90}-S_{120}.}
 \label{eq:l9s102:barbero-peso-generado}
\end{equation}
Thus, the weight \(12:-1\) is the image of the fundamental turn: twelve sectors and
one oriented boundary. It does not arise from imposing \(1/24\) beforehand, from a
decimal comparison or from the conventional Barbero--Immirzi value.

\Needspace{12\baselineskip}
\begin{theorem}[Dodecaphase determination of the Barbero--Immirzi functional]
\label{thm:l9s102:barbero-funcional}
The fundamental chain \eqref{eq:l9s102:barbero-cadena-dodecafasica} determines
the oriented pair \((a,b)=(12,1)\). The coefficient of \((1-q)^{-1}\) in its image is
then \(1/24\), as an exact output of the already constructed functional. Therefore,
after evaluating the conjugate channels \(q_\pm\), the dimensionless output is
\begin{equation}
 \boxed{
 \gamma_{\rm HMT}=\Gamma_{6\to8}
 =\frac{\pi_{\rm HMT}AC^*}{180}
   +12\Delta S_{90}-\Delta S_{120}.}
 \label{eq:l9s102:barbero-gamma}
\end{equation}
This formula receives the already generated angular pair and does not use
 the conventional value of the Barbero--Immirzi parameter as input.
\end{theorem}

\begin{proof}
The sum of \eqref{eq:l9s102:barbero-cadena-dodecafasica} contains exactly
the twelve sectors of \(C_{12}\), while the return seam appears
only once and with negative orientation. The reader
\eqref{eq:l9s102:barbero-lector-retorno} therefore produces
\eqref{eq:l9s102:barbero-peso-generado}, before any pole calculation. For
each \(m>0\),
\[
 p_1(S_m)
 =\lim_{q\to1}(1-q)\frac{q^m}{1-q^{3m}}
 =\frac1{3m}.
\]
Consequently,
\[
 p_1\bigl(\mathscr R(c_{12}^{+})\bigr)
 =\frac{12}{270}-\frac1{360}
 =\frac{48-3}{1080}=\frac1{24}.
\]
The coefficient \(p_1\) is referred to \((1-q)^{-1}\).
In the complex coordinate \(q-1\),
\[
 \operatorname*{Res}_{q=1}S_m(q)=-\frac1{3m},
 \qquad
 \operatorname*{Res}_{q=1}\mathscr R(c_{12}^{+})(q)=-\frac1{24}.
\]
The subsequent Diophantine equation
\[
 \frac a{270}-\frac b{360}=\frac1{24}
 \quad\Longleftrightarrow\quad 4a-3b=45
\]
has the positive integer solutions
\((a,b)=(12+3t,1+4t)\), \(t\ge0\); the unique minimal solution is
\((12,1)\). This count arithmetically certifies the pair that the fundamental
orbit had already generated: any \(t>0\) alters the multiplicity of
the sectors or the seam and ceases to represent one TPK turn.
Finally, substituting \(q_\pm\) into the oriented image and adding the previously generated
angular component proves \eqref{eq:l9s102:barbero-gamma}.
\end{proof}

The causal sequence is fixed without exchanging roles:
\begin{equation}
\begin{aligned}
 (\mathcal F_6\hookrightarrow\mathcal F_8,C_{12},\partial_{\rm ret})
 &\longrightarrow (90,120,-1/12,c_{12}^{+})\\
 &\longrightarrow \mathscr R(c_{12}^{+})=12S_{90}-S_{120},\\
 \mathscr R(c_{12}^{+})
 &\longrightarrow p_1\bigl(\mathscr R(c_{12}^{+})\bigr)=\frac1{24},\\
 \bigl(\mathscr R(c_{12}^{+}),A,C^*,q_\pm\bigr)
 &\longrightarrow \Gamma_{6\to8}=\gamma_{\rm HMT}
 \longrightarrow \widehat{\mathcal A}_{\rm phys}.
\end{aligned}
 \label{eq:l9s102:barbero-cadena-completa}
\end{equation}
The angular pair and its channels come from the preceding construction.
Evaluation of the pole coefficient is a check
of the return functional; the conjugate evaluation uses the
previously generated angular coordinates.
The subsequent realization in the Holst action recognizes the same dimensionless
section; it does not intervene in its selection.

For the two boundary sectors, let
\begin{equation}
 N_m=\sum_{e\in\partial A_m}N_e,
 \qquad m\in\{90,120\},
 \qquad \operatorname{Spec}(N_e)=\{0,1,2\}.
 \label{eq:l9s102:n90-n120}
\end{equation}
The physical area operator is written
\begin{equation}
 \boxed{
 \frac{\widehat{\mathcal A}_{\mathrm{phys}}}{\ell_P^2}
 =\gamma_{\mathrm{HMT}}a_0(N_{90}+N_{120}).}
 \label{eq:l9s102:area-operador}
\end{equation}
On \(36\) links qutrit,
\begin{equation}
 \operatorname{Spec}(N_{90}+N_{120})=\{0,1,\ldots,72\},
 \label{eq:l9s102:area-numero-spectrum}
\end{equation}
and the multiplicity of the eigenvalue \(n\) is the coefficient of \(z^n\) in
\((1+z+z^2)^{36}\). From here it follows
\begin{equation}
 \operatorname{Spec}\!\left(
 \frac{\widehat{\mathcal A}_{\mathrm{phys}}}{\ell_P^2}\right)
 =\{n\gamma_{\mathrm{HMT}}a_0:0\le n\le72\}.
 \label{eq:l9s102:area-spectrum}
\end{equation}

\section{Total Geometry and Origin Memory}
\label{sec:l9s102:area-hamiltoniano}

The modular Hamiltonian of the boundary is
\begin{equation}
 \boxed{
 K_A=-\log\rho_A
 =cI+\Delta_{\mathrm{mod}}(90N_{90}+120N_{120}).}
 \label{eq:l9s102:hamiltoniano-modular}
\end{equation}
Let
\begin{equation}
 N=N_{90}+N_{120},
 \qquad
 D=90N_{90}+120N_{120}.
 \label{eq:l9s102:n-d}
\end{equation}
The transition matrix
\(\left(\begin{smallmatrix}1&1\\90&120\end{smallmatrix}\right)\)
has determinant \(30\), and its inverse reconstructs the sectors:
\begin{equation}
 \boxed{
 N_{90}=4N-\frac{D}{30},
 \qquad
 N_{120}=\frac{D}{30}-3N.}
 \label{eq:l9s102:reconstruccion-sectores}
\end{equation}
The area publishes the total number of excitations; the modular Hamiltonian
publishes their provenance sector. The joint pair reconstructs the complete
occupation of the boundary.

\section{Bilayer state and oriented information}
\label{sec:l9s102:estado-bicapa}

Let \(R=R^*=R^{-1}\) be the sheet involution and
\(y=C^*_{\mathrm{rad}}\). The normalized state is
\begin{equation}
 \boxed{
 \rho_y=\frac{e^{-yR}}{2\cosh y}.}
 \label{eq:l9s102:rho-y}
\end{equation}
Its modular Hamiltonian and its entropy are
\begin{equation}
 K_y=-\log\rho_y=\log(2\cosh y)I+yR,
 \label{eq:l9s102:ky}
\end{equation}
\begin{equation}
 S(\rho_y)=\log(2\cosh y)-y\tanh y.
 \label{eq:l9s102:entropia-y}
\end{equation}
The inversion \(y\mapsto-y\) leaves entropy fixed and produces the oriented
informational distance
\begin{equation}
 \boxed{
 D(\rho_y\Vert\rho_{-y})=2y\tanh y.}
 \label{eq:l9s102:entropia-relativa-hojas}
\end{equation}
The even component measures the spectral distribution; the odd component measures
the informational cost of reversing the orientation.

\section{How much informational energy and gravitational cell of medium action}
\label{sec:l9s102:horizonte-media-accion}

The periodicity of \(108\) states determines
\begin{equation}
 t_P=\frac{54}{\pi}t_0,
 \qquad
 E_P=\frac{\hbar}{t_P}=\frac{h}{108t_0}.
 \label{eq:l9s102:ep-tp}
\end{equation}
The thermal realization of the tritico clock satisfies
\begin{equation}
 \boxed{
 E_P=k_B\Theta_{\mathrm{clk}}\log3.}
 \label{eq:l9s102:planck-trit}
\end{equation}
For the horizon of radius \(R\),
\begin{equation}
 E_\Lambda(R)=\frac{c^4R}{2G},
 \qquad
 E_\Lambda(R)=T_H(R)S_H(R).
 \label{eq:l9s102:energia-horizonte}
\end{equation}
In the Planck cell,
\begin{equation}
 \boxed{
 E_\Lambda
 =T_HS_H
 =\frac{E_P}{2}
 =\frac{k_B\Theta_{\mathrm{clk}}\log3}{2}.}
 \label{eq:l9s102:confluencia-horizonte}
\end{equation}
Therefore,
\begin{equation}
 \boxed{
 E_\Lambda t_P=\frac{\hbar}{2},
 \qquad
 E_\Lambda T_{108}=\frac h2.}
 \label{eq:l9s102:media-accion}
\end{equation}
Vacuum energy, area entropy, chronological temperature, and
action share a single confluence cell with typed domains.

When \(S_H/k_B=\pi\), the multiplicity efectiva satisfies
\begin{equation}
 \Omega_{\mathrm{eff}}=e^{S_H/k_B}=e^\pi.
 \label{eq:l9s102:omega-e-pi}
\end{equation}
The hyperbolic Euler operator possesses
\begin{equation}
 \operatorname{Spec}(e^{\pi H})=\{e^\pi,e^{-\pi}\}.
 \label{eq:l9s102:e-pi-spectrum}
\end{equation}
Equality of the scalars establishes spectral confluence; the
intertwiner between microscopic horizon multiplicity and Perron
dynamics belongs to its specific physical residence.
